\section{Introduction}\label{intro}

GPUs are in high demand for large-scale Machine Learning (ML). As ML models continue to grow exponentially in size, they require an increasing number of GPUs to train and fine-tune. This high demand makes it difficult for model developers to allocate the desired number of accelerators to train models at high throughput in public clouds or enterprise clusters~\cite{Um24Metis, strati24crossregion, Zongheng23Skypilot}. 
Datacenters typically host a variety of GPU types and generations, spread across geographic regions~\cite{gcpgpus, Jia2022Whale, weng22mlaas, Jeon19philly, wu2024falcon}. Yet model developers tend to restrict model training to \textit{homogeneous} clusters of GPUs, since state-of-the-art distributed training frameworks like Megatron-LM~\cite{shoeybi2020megatronlm} and DeepSpeed~\cite{megatron-deepspeed} assume homogeneous GPUs and inter-node bandwidth. The demand for large, homogeneous GPU clusters compounds the scarcity of high-end GPUs.

\begin{figure}[t]
\centering
  \includegraphics[width=\linewidth]{figures/motivation/intro/OPT-350.pdf}\label{fig:motivation_het}%
\vspace{-5pt}
\caption{When homogeneous resources are limited, using heterogeneous GPUs (A100, V100) or multiple zones can increase OPT-350M training throughput (c3-c4) at low cost (black line). However, when the resource topology and job parallelization are not well selected, iteration time and monetary cost may increase significantly (c5-c6).} 
\vspace{-3pt}
\label{fig:motivation_het}
\end{figure}

Allowing a training job to run on heterogeneous GPU types and/or GPUs distributed across zones (i.e., with heterogeneous inter-node bandwidth) can give model developers access to \textit{more GPUs} per job to increase training throughput. 
For example, consider a model developer who seeks to maximize training throughput by using 32 A100 GPUs (\textcolor{blue}{\texttt{c2}} in Figure~\ref{fig:motivation_het}), but discovers that only 16 A100s (\textcolor{blue}{\texttt{c0}} in Figure~\ref{fig:motivation_het}) are currently available in one zone. Using heterogeneous GPU generations (e.g., \textcolor{teal}{\texttt{c3}} uses an additional 16 V100s in the zone) or multi-zone configurations (e.g., \textcolor{teal}{\texttt{c4}} uses 32 A100s spread across 2 zones within a region) increases throughput by 1.15$\times$ and 1.87$\times$, respectively, with a moderate increase in monetary cost per iteration (black line). 





However, supporting heterogeneous, geo-distributed resources introduces several challenges.
First, heterogeneous GPU types and placements across zones exponentially expand the configuration search space. 
Searching for an optimal configuration requires jointly optimizing the \textit{resource allocation} and \textit{job parallelization plan} (e.g., data/pipeline/tensor parallelism degrees). For example, in Figure~\ref{fig:motivation_het}, although \textcolor{red}{\texttt{c5}} uses the same number of GPUs as \textcolor{teal}{\texttt{c3}}, it achieves much lower throughput due to a suboptimal parallelization plan. 
The search must also consider the cost caused by extra resources and data transfers. Cloud providers charge significant fees for inter-zone and inter-region data transfers~\cite{aws-interzone, gcp-interzone, az-interzone}, which impacts geo-distributed configurations. In Figure~\ref{fig:motivation_het}, \textcolor{red}{\texttt{c6}} uses the same number of GPUs and parallelization strategy as \textcolor{teal}{\texttt{c4}}, but spreads training across regions (instead of zones within the same region), which increases cost. 

Furthermore, as resource availability changes dynamically in datacenters, due to variable demand and node failures or preemptions, it is necessary to navigate this vast search space quickly ~\cite{Thorpe23Bamboo, wagenlander24tenplex}. Figure~\ref{fig:trace} shows the varying number of A100 GPUs that we were able to allocate over an 8-hour period, out of the 8 A100s that we continuously requested in 2 different zones in Google Cloud. Such changes in resource availability require frequently re-evaluating the job configuration search space. Thus, model developers require a system to quickly navigate the large search space of heterogeneous (and homogeneous) configurations with dynamic resource availability, optimizing for the user's performance-cost objective, while satisfying constraints (e.g., budget limits).  


Second, profiling many candidate configurations to evaluate their throughput is prohibitively expensive and time-consuming. Hence, it is critical to accurately estimate the iteration time of a candidate configuration. This is challenging in heterogeneous environments, as differences in peak FLOPS, memory capacity, CPU-GPU interconnects, number of GPUs per node, and inter-node bandwidth typically lead to \textit{stragglers}, which can significantly limit throughput. 
More importantly, variability in memory capacity per GPU may cause \textit{out-of-memory (OOM) errors} in some GPUs, disrupting the entire training job. Simulating iteration time and checking that job configurations are valid (i.e., will not cause OOM) requires correctly modeling stragglers and per-GPU memory footprints.



Finally, after finding an appropriate resource allocation and parallelization plan for a training job given the available resources, model developers need to be able to run this job configuration in a distributed training framework (such as Megatron~\cite{shoeybi2020megatronlm}). We find that optimal configurations for jobs running on heterogeneous resource topologies often include heterogeneous parallelism degrees per stage to load-balance the compute and memory capacity of different GPU types. 
Today's state-of-the-art distributed training frameworks need to be adapted to support such heterogeneous job configurations. Furthermore, as resource availability can change frequently (e.g., when using spot instances~\cite{Thorpe23Bamboo}), the training framework must be able to quickly reconfigure jobs.


Existing systems do not adequately solve these challenges. First,   
current works do not co-optimize the resource allocation with the job parallelization plan. Instead, systems like Aceso~\cite{Liu24Aceso}, Galvatron~\cite{Miao22Galvatron}, and others in Table~\ref{table:planners} expect the user to select a fixed resource allocation for which the system recommends a job parallelization plan. Most systems also do not consider heterogeneous resource topologies. 
Recent systems like Atlas~\cite{palak2024atlas}, DTFM~\cite{yuan2023decentralized}, Metis~\cite{Um24Metis}, and FlashFlex~\cite{yan2024flashflex} optimize parallelization for heterogeneous GPUs or geo-distributed setups, but they suffer from prohibitively long search times (up to hours for configurations with 10s of GPUs)~\cite{Um24Metis}, or suboptimal cost functions~\cite{yan2024flashflex, yuan2023decentralized}, making them unsuitable for environments with dynamic resource availability. 
Second, existing systems rely on inaccurate simulators to estimate the training throughput and memory footprint of candidate configurations. For example, Varuna~\cite{Athlur22Varuna} overlooks significant memory sources (e.g., memory needed by the optimizer, communication, etc) when estimating memory footprint, hence recommending configurations that lead to OOM errors. Finally, state-of-the-art distributed training frameworks like Megatron-LM~\cite{shoeybi2020megatronlm} are slow to reconfigure jobs and do not support heterogeneous job parallelization plans or different microbatch sizes per GPU, which is necessary to maximize throughput and minimize cost in heterogeneous clusters.

To this end, we propose \textit{\sailor}, a system for efficient large-scale training over heterogeneous resources with dynamic availability. \sailor\footnote{\sailor is available at \href{https://github.com/eth-easl/sailor}{https://github.com/eth-easl/sailor}} consists of three components: a configuration planner, a simulator, and a distributed training framework. The \sailor planner navigates the search space of resource allocations and job parallelization plan combinations. It recommends configurations that optimize a user-defined objective (e.g., max throughput or min cost) under constraints (e.g., max budget or min throughput). The planner considers heterogeneous GPU and machine types and geo-distributed setups. The planner uses the simulator to accurately model iteration time and memory footprint for any given configuration. 
Through a combination of dynamic programming and search space pruning with effective heuristics, the planner finds solutions within seconds even for allocations with 100s of GPUs and varying degrees of heterogeneity. This allows \sailor to quickly adapt plans based on resource availability. Finally, the \sailor training framework adds support for heterogeneous configurations to execute the planner’s configurations. It also adds support for fault tolerance and elasticity, enabling adaptation to changes in resource availability. Together, these components enable \sailor to efficiently automate large-scale training in homogeneous, heterogeneous, and/or dynamic resource environments.

We evaluate \sailor in various setups and compare it extensively to prior works. 
To the best of our knowledge, our work is the first to compare the major open-source ML training planners proposed to-date (Table \ref{table:planners}) in homogeneous and heterogeneous scenarios. We show that \sailor can find resource allocations and job parallelization plans that result in higher throughput than baselines in both homogeneous and heterogeneous clusters.
We show how \sailor can leverage heterogeneous resources to improve throughput by 1.1-2.87$\times$ compared to the heterogeneity-aware baselines (Metis, FlashFlex, AMP), while maintaining search times of 10s of seconds compared to minutes or hours needed by the baselines. We also show \sailor's ability to increase performance and reduce cost in geo-distributed setups compared to DTFM by 5.9$\times$ and 9.8$\times$, respectively. Finally, we demonstrate \sailor's ability to minimize monetary cost given throughput constraints, resulting in 40\% cost savings compared to the second-best-performing baseline.